
\chapter{Textual Data}\label{sec-textual-data}

\textbf{Learning Objectives}

\emph{Overview starts}

\begin{itemize}
\tightlist
\item
  Understanding the fundamental concepts of representing, manipulating, and
  working with text, such as data structures (vectors, arrays, and data
  frames), text strings, and characters.
\item
  Learning how to import a text dataset and display its strings, and handle
  common issues in UGC, such as misspelling, typos, and non-standard
  formats.
\item
  Acquiring the skills to obtain and analyze information about strings within
  a dataset.
\item
  Developing an understanding of regular expressions in R and their
  practical utility.
\item
  Practicing by using regular expression tools including quantifiers,
  boundaries, operators, character classes, and sequences.
\end{itemize}

\emph{Overview ends}

\section{Text manipulation}\label{text-manipulation}

\subsection{Introduction}\label{introduction}

It is rare for a marketing analyst to work with perfectly crafted texts.
Since the focus of this book is on analyzing user-generated content
(UGC), such as online reviews of products or services, blog entries, or
social media posts, from a marketing perspective, we cannot expect any of
these texts to be free from linguistic errors, typos, or misspellings.
Unless they are revised or edited by professionals, these texts are
likely to be imperfect (Moens, Li, and Chua 2014).

Such linguistic errors in UCG can impact the analyses of
these texts and, in turn, the marketing implications derived from them negatively.
For instance, typos in reviews can shift the overall sentiment. Mistakes
can lead to the misinterpretation of texts and potentially omit essential
product information. Different spellings of the product name or features
can make product comparisons inconsistent. Incorrectly spelled keywords
might make it challenging for companies to track conversations and
analyze customer sentiment (Berger et al. 2020). In some contexts,
misspelling can even be used as an additional measure
to reveal information about the communicator (Netzer, Lemaire, and
Herzenstein 2019). These are not the only problems found in UGC.
There are also punctuation errors, incorrect capitalization, nonstandard
abbreviations and acronyms, misused emojis, homophones, slang, jargon,
incomplete sentences, and many more elements which can be misinterpreted. All in all, there are countless ways to write a text which present problems for textual analysis.

Thus, we often encounter situations when we need to search for, select,
and partly modify or transform words, numbers, or whole phrases. This
process, known as text manipulation, requires specific skills. It is
necessary when aggregating text from different sources, such as from
online product reviews with authors from different backgrounds,
cultures, or with varying styles. The main goal of altering text in reviews is
to create documents that are error-free, uniform, and ready to be
analyzed. This is important because humans can correct
misspelled words mentally but are slow, computers cannot make such a distinction
but are fast. If the text is not polished, the resulting statistics will
be flawed: for example, different spellings of the same word will be
treated as separate words by the computer.

To catch common spelling errors, most of us use a spell
checker of some form when we compose a text document, or at least we
should have it activated. A spell checker is a function built into
common text writing programs that detect misspellings by comparing the
words in a drafted document with a database containing standard
spellings. Among others, a popular text-mining package that includes
prepackaged spellers to correct mistakes is the Enchant speller
(AbiWord Project 2024). Modern applications use artificial intelligence
algorithms to do this, but although using such automatic tools for large-scale
correction may seem convenient, it comes with risks. The main
challenge is that these tools may only sometimes determine context. This
can lead to incorrect changes and generate unwanted errors. With
advancements in AI, though, contextual understanding is becoming more accurate, thereby reducing such risks (Hu et al. 2020).

The discrepancy between English language variants that have a disparity
in standardization, such as British English (GB) or American English
(US) spelling, is another problem (Hitchings 2008). When dealing
with aggregated data, the probability of encountering both in the same
corpus is high. Thus, a method to make them uniform is desirable. We can create
ad-hoc functions or procedures to manipulate text to overcome these
difficulties. The operations we employ to achieve take three forms: additive, subtractive, and substitutive. In the first case,
we can add characters, strings, or numbers to a specific string
anywhere, such as at the beginning, in the middle or at the end of a word, and before or
after a specific character. In the same way, we can subtract or
substitute them to obtain the desired result. In essence,
 \textbf{text is your data}, and managing it effectively is
essential for accurate analysis.

Now we will briefly introduce some basic concepts of how (textual) data are
stored inside R. Readers who are merely interested in applications can
skip this section. Later, we will delve into more advanced techniques
and functions to help us manipulate textual data effectively.

\subsection{Some basics on measurement representation in
R}\label{some-basics-on-measurement-representations-in-r}

Data is more than knowledge (Tuomi 1999). The Oxford dictionary defines
\textbf{data} (Coombs 1964) as: ``The quantities, characters, or symbols
on which operations are performed by a computer, which may be stored and
transmitted in the form of electrical signals and recorded on magnetic,
optical, or mechanical recording media.'' It is easy to see the
\textbf{quantities} as numbers, the \textbf{characters} as letters, and
the \textbf{symbols} as images, and consider all these as \textbf{types of data}. In computer programming, we deal with two fundamental entities for
accessing stored data: \textbf{variables} and \textbf{constants} (Ihaka
and Gentleman 1996). Variables in programming are used to point to
stored data by assigning them an understandable, meaningful, and concise
name. For example, we could use the variable name \texttt{temperature} to refer
to a sequence of temperature values that you have measured. Then, the
variable can be referenced and modified as needed by other computer
routines. The \textbf{constant} represents a fixed value that cannot be
altered but can be used in other functions. A classic example is
\textbf{pi}, the ratio of the circumference of a circle to its diameter,
commonly rounded down to \textbf{3.14}. Data need to be structured
(Arasu and Garcia-Molina 2003) to facilitate easy manipulation. For the
sake of precision, you need to understand that a \textbf{variable}
represents a reference to a structure which has been stored; it does not store the
data, it just gives them a name.

\subsubsection{Vectors}\label{vectors}

A \textbf{vector} is the simplest form of data structure in R, commonly
used alongside lists, matrices, arrays, and data frames. A vector can be
\textbf{atomic}, meaning that it contains elements that are all of the
same data type, for example logical, integer, double, or character data. It
can contain a single value, or it can hold multiple values. In
Figure~\ref{fig-array}, we can make an analogy between a vector, which
contains variables of a single data type, and a single-column chest of
drawers filled with socks. Each drawer represents an element in the
vector, and we have named this vector \texttt{drawers}.

To populate our vector, we assign values to it using the variable name
(\texttt{drawers}) and a discrete number of differently colored socks:
\textbf{23 red}, \textbf{32 blue}, \textbf{15 green}, \textbf{41
yellow}, and \textbf{12 white}. We use the function, \texttt{c()}, which
stands for \textbf{combine}, to group these ``socks'' together. The
assignment is done using the \texttt{\textless{}-} operator:
\texttt{drawers\ \textless{}-\ c(23,\ 32,\ 15,\ 41,\ 12)}.

To access the content of any drawer, we simply use the variable name and
specify the desired element(s) using \textbf{square brackets}. For
instance, \texttt{drawers{[}3{]}} retrieves the content of the third
element of the vector variable, which is \textbf{15 green socks}.
Remember, what appears in the R console are numbers, not actual socks.
If we want to know the contents of the second and fifth drawers, we
type: \texttt{drawers{[}c(2,\ 5){]}}, which will return the number of
socks in both the second and fifth elements of the vector. As you can
see, our vector contains \textbf{only socks}, not shirts or pants.
This shows that a vector can contain only elements of the same type
(e.g., all numbers, all integers, or all characters). On top of that, it
is \textbf{one-dimensional}, because it is flat and linear.

\subsubsection{Arrays}\label{arrays}

Let us introduce another piece of terminology: the \textbf{array}. In R, an
\textbf{array} is like a vector but can have more than one dimension. An example of a two-dimensional array is a \textbf{matrix}.
Technically, they are the same thing, but R likes to keep these two
concepts separate due to the significant role matrices play in
statistical computation. Obviously, we can have arrays with more than
two dimensions, but visualizing these higher-dimensional structures can
be challenging, if not impossible.

\begin{figure}

\centering{

\includegraphics[width=0.99\textwidth,height=\textheight]{figure/chapter2/array.png}

}

\caption{\label{fig-array}An array of a chest of drawers
\newline Authors' own figure}

\end{figure}%

In R, a \textbf{list} is a data structure that can store elements of varying types—such as numbers, strings, \textbf{vectors}, or even other lists. The key difference between a \textbf{list} and a \textbf{vector}
lies in the types of data they can contain. While a \textbf{vector} must
consist of elements of the \textbf{same data type}, a list can hold
\textbf{heterogeneous data types}, including vectors, other lists, and
even lists of lists. This flexibility makes lists a versatile data
structure in R.

\subsubsection{Vector Types in R}\label{vector-types-in-r}

R language categorizes data into six basic atomic vector types: three
numerical types (integer, numeric, complex), a logical type (true or
false), a character type, and a raw type. For a more comprehensive
understanding of these types, refer to the works of Kabacoff (2022) and
Wickham (2019), which provide detailed explanations and examples.

These types are named intuitively to reflect the kind of data
they hold. For instance, \textbf{integer vectors} store whole numbers,
\textbf{numeric vectors} can hold decimal numbers, \textbf{complex
vectors} are for complex numbers, \textbf{logical vectors} contain
Boolean values (TRUE or FALSE), \textbf{character vectors} store text,
and \textbf{raw vectors} are used for raw bytes. As an analogy, think of
the data type as the material from which each chest of drawers is made.
For instance, some might be made of wood (representing integers), others
of plastic (representing characters), and some others of metal
(representing logical values like TRUE or FALSE).

\subsubsection{Data Frames}\label{data-frames}

A \textbf{data frame} in R is structured like a table, and is similar to a table in many standard spreadsheet programs
such as Microsoft Excel, Google Sheets, Apple Numbers, or LibreOffice
Calc.

Just as a table in these programs is composed of rows and columns, an R
\textbf{data frame} is structured in the same way. Each \textbf{column} in a
data frame can be considered to be a vector, holding data of the same type,
much like a spreadsheet column containing numbers, text, or dates. The
\textbf{rows} in a data frame correspond to individual records or
observations, similar to how each row in a spreadsheet represents a
specific item, person, or event.

To continue with our analogy between a data frame in R and a chest of
drawers in the physical world, let us consider each row in the data
frame to represent a person, see Figure~\ref{fig-dataframe}. In
technical terms, the data in each \textbf{row} is called an
\textbf{observation}. A \textbf{column} represents a specific feature,
called an \textbf{attribute}. Each person, or row, is linked to various features,
which in our analogy are items like socks, t-shirts, and pants. For
instance, Aida might have 23 socks, 16 t-shirts, and 8 pairs of pants.

\subsubsection{Manipulating Data Frames}\label{manipulating-data-frames}

In this context, the chest of drawers can be thought of as our data
frame in R. Each drawer or column represents a different item of
clothing. We access these items in R in a similar way to opening a specific
drawer of the piece of furniture. For example, if we want to check how many t-shirts
Carl has, and Carl is the third person in our data frame, we would use
the command \texttt{chest{[}3,\ 2{]}} in R. This command is akin to
opening the second drawer (t-shirts) of the third row (Carl) of our
chest of drawers.

We can also `slice' this data frame, which is just like selecting specific
drawers from our chest of drawers. If we want to see only the t-shirts and pants
for everyone, we use \texttt{chest{[},\ c(2,\ 3){]}}. The empty space
before the comma indicates that we are selecting all the rows or, in our
analogy, all the people. After the comma, the \texttt{c(2,\ 3)} specifies
that we only want to look at the second and third columns, which are
t-shirts and pants.

In this analogy, instead of actual socks, t-shirts, and pants, we are
dealing with different types of data in our data frame. So, where you might
imagine long socks in a drawer, in R, this could correspond to a column
of strings or numbers representing that particular item or attribute.

\begin{figure}

\centering{

\includegraphics[width=0.99\textwidth,height=\textheight]{figure/chapter2/dataframe.png}

}

\caption{\label{fig-dataframe}A data frame containing various items
\newline Authors' own figure}

\end{figure}%

Imagine that your chest of drawers can be part of different sets of furniture. One chest might belong to a set for a bedroom (similar to a data frame
\textbf{class}), another might be part of a set of office furniture
(similar to a list \textbf{class}), and so on. The set of furniture
(\textbf{class}) it belongs to determines how the chest of drawers is
used. A bedroom chest of drawers might be used to store clothes, while an
office version might be organized to store files and documents. In R,
the \textbf{class} of an object determines how it behaves with different
functions and operations. For example, functions treat a data frame
(a bedroom set) differently than a basic vector (a standalone chest of
drawers). The difference between \textbf{data type} and \textbf{class}
is the following. The first refers to the \textbf{low-level internal
structure} of how R stores data in memory. A variable can be checked for
its type using the \texttt{typeof()} function. while the second is a
\textbf{high-level} \textbf{abstraction} that focuses on how an object
behaves in R. Because the word \textbf{class} is related to object
oriented programming, it defines how functions and methods interact with
an object. As a reminder, an \textbf{object} is an element which stores
information, it is referenced by a name, which is the name of the variable.

\subsubsection{Empty Strings and Character
Vectors}\label{empty-strings-and-character-vectors}

Just as we use a drawer to store garments, we can use specific
structures in programming to store text strings. Let us imagine we have
a chest of drawers, but we will leave these drawers empty. In the
context of programming, particularly in R, text strings are categorized
as belonging to a \textbf{character class}. So we use another data
type now, a \textbf{character}, not an integer.

Let us consider an \textbf{empty string}, which is similar to having a chest of
drawers with only one column of drawers, and those drawers are empty. We
can assign an \textbf{empty string} to a variable; let us call it
\texttt{empty\_string}. This would be shown in R as
\texttt{empty\_string\ \textless{}-\ ""}, signifying a string that
contains no characters. Paradoxically, even though the string is empty,
it belongs to the \textbf{character class} in R.

When we are uncertain about the class type of a variable, we can easily
check it using the \texttt{class()} function. For instance, typing
\texttt{class(empty\_string)} would confirm that empty\_string is indeed
a \textbf{character class} despite being empty. This is similar to
confirming that an empty drawer is still a part of the chest of drawers,
designed to hold items even when it currently does not.

\emph{Program code starts}

\begin{Shaded}
\begin{Highlighting}[]
\CommentTok{\# Assign an empty string to the variable \textquotesingle{}empty\_string\textquotesingle{}}
\NormalTok{empty\_string }\OtherTok{\textless{}{-}} \StringTok{""}

\CommentTok{\# Display the value of \textquotesingle{}empty\_string\textquotesingle{}}
\NormalTok{empty\_string}
\end{Highlighting}
\end{Shaded}

\begin{verbatim}
[1] ""
\end{verbatim}

\begin{Shaded}
\begin{Highlighting}[]
\CommentTok{\# Check and display the class of \textquotesingle{}empty\_string\textquotesingle{}}
\FunctionTok{class}\NormalTok{(empty\_string)}
\end{Highlighting}
\end{Shaded}

\begin{verbatim}
[1] "character"
\end{verbatim}

\emph{Program code ends}

Similarly, we can create an empty character vector in R using the
\texttt{character()} function by setting its length argument to zero. It
is important to note that \texttt{empty\_character} is just a name we
have assigned to our variable (object) for clarity and understanding. We
could have named it anything, such as \texttt{john\_doe},
\texttt{ala\_bama}, or \texttt{bugs\_bunny}. Choosing a descriptive name
is beneficial. When we revisit the code later and need an immediate
glimpse of what a specific variable might contain, the appropriate naming of
variables makes them easier to read and understand.

\emph{Program code starts}

\begin{Shaded}
\begin{Highlighting}[]
\CommentTok{\# Create an empty character vector with a length of zero}
\NormalTok{empty\_character }\OtherTok{\textless{}{-}} \FunctionTok{character}\NormalTok{(}\AttributeTok{length =} \DecValTok{0}\NormalTok{)}

\CommentTok{\# Display the content of \textquotesingle{}empty\_character\textquotesingle{}}
\NormalTok{empty\_character}
\end{Highlighting}
\end{Shaded}

\begin{verbatim}
character(0)
\end{verbatim}

\begin{Shaded}
\begin{Highlighting}[]
\CommentTok{\# Check and display the class of \textquotesingle{}empty\_character\textquotesingle{}}
\FunctionTok{class}\NormalTok{(empty\_character)}
\end{Highlighting}
\end{Shaded}

\begin{verbatim}
[1] "character"
\end{verbatim}

\emph{Program code ends}

What if we want to create a character vector in R that is initialized
(meaning it occupies spots but contains only empty placeholder values)?
The procedure can be thought of as preparing a chest of drawers, where
each drawer is meant to hold a character string, but initially, all the
drawers are empty. In programming terms, we are talking about
initializing an empty character vector and we can set it to be a specific
length, for example to have ten elements, precisely like a chest of drawers with
ten drawers. At this stage, each `drawer' (or element in the vector)
contains an \textbf{empty string} \texttt{("")} -- ready to hold data,
but currently contains no values.

\emph{Program code starts}

\begin{Shaded}
\begin{Highlighting}[]
\CommentTok{\# Create a character vector with 10 elements, initialized as empty strings}
\NormalTok{character\_vector }\OtherTok{\textless{}{-}} \FunctionTok{character}\NormalTok{(}\DecValTok{10}\NormalTok{)}

\CommentTok{\# Display the contents of \textquotesingle{}character\_vector\textquotesingle{}}
\NormalTok{character\_vector}
\end{Highlighting}
\end{Shaded}

\begin{verbatim}
 [1] "" "" "" "" "" "" "" "" "" ""
\end{verbatim}

\emph{Program code ends}

\subsubsection{Combining and Indexing
Strings}\label{combining-and-indexing-strings}

We can populate our newly created character vector with words or strings
by assigning them to specific positions within the vector. This is done
through indexing, which allows us to access and modify elements at
particular locations in the vector. For example, if we want to place the
word `five' in the fifth position of our vector, we use indexing. This is done by adding square brackets \texttt{{[}{]}} after
the vector's name, \texttt{character\_vector}, and specifying the position we want with a number. In this case, we use the number 5 to
indicate the fifth position. The corresponding R code would be
\texttt{character\_vector{[}5{]}\ \textless{}-\ \textquotesingle{}five,\textquotesingle{}}
which assigns the string `five' to the fifth element of the vector.

\emph{Program code starts}

\begin{Shaded}
\begin{Highlighting}[]
\CommentTok{\# Assign the string "fifth" to the fifth element of the \textquotesingle{}character\_vector\textquotesingle{}}
\NormalTok{character\_vector[}\DecValTok{5}\NormalTok{] }\OtherTok{\textless{}{-}} \StringTok{"fifth"}

\CommentTok{\# Display the contents of \textquotesingle{}character\_vector\textquotesingle{}}
\NormalTok{character\_vector}
\end{Highlighting}
\end{Shaded}

\begin{verbatim}
 [1] ""      ""      ""      ""      "fifth" ""      ""      ""      ""     
[10] ""     
\end{verbatim}

\emph{Program code ends}

As previously mentioned, a vector in R is designed for homogeneous
storage, meaning it can only hold one type of data. For instance,
if we add a string to a numerical vector, R will automatically convert
all numerical values in the vector to strings. This is because a vector
must maintain a consistent data type across all its elements, or
`drawers', as we might metaphorically describe them. We can
convert vector types manually if we require precise control over our types of data.
To check if a vector is a character type, we use the function
\texttt{is.character()}. Moving from one type to another, such as from
numerical to character, is achieved using functions like
\texttt{as.character()}. This operation might be performed when we want
to transform numerical value types into character types. While number types
are used for computations, numbers as text (character type) are
interpreted as words.

Consider the following example where two variables, \texttt{x} and
\texttt{y}, hold different types of data: character and numerical,
respectively. We can check if they are a character
type, examine their class, convert one type to another, and then verify
if the conversion was successful.

\emph{Program code starts}

\begin{Shaded}
\begin{Highlighting}[]
\CommentTok{\# Assign a string to variable \textquotesingle{}x\textquotesingle{}}
\NormalTok{x }\OtherTok{\textless{}{-}} \StringTok{"really, is this a character?"}

\CommentTok{\# Assign the result of the sum 9+8 to variable \textquotesingle{}y\textquotesingle{}}
\NormalTok{y }\OtherTok{\textless{}{-}} \DecValTok{9} \SpecialCharTok{+} \DecValTok{8}

\CommentTok{\# Check if \textquotesingle{}x\textquotesingle{} is of character type}
\FunctionTok{is.character}\NormalTok{(x)}
\end{Highlighting}
\end{Shaded}

\begin{verbatim}
[1] TRUE
\end{verbatim}

\begin{Shaded}
\begin{Highlighting}[]
\CommentTok{\# Check if \textquotesingle{}y\textquotesingle{} is of character type}
\FunctionTok{is.character}\NormalTok{(y)}
\end{Highlighting}
\end{Shaded}

\begin{verbatim}
[1] FALSE
\end{verbatim}

\begin{Shaded}
\begin{Highlighting}[]
\CommentTok{\# Get the class of \textquotesingle{}x\textquotesingle{}}
\FunctionTok{class}\NormalTok{(x)}
\end{Highlighting}
\end{Shaded}

\begin{verbatim}
[1] "character"
\end{verbatim}

\begin{Shaded}
\begin{Highlighting}[]
\CommentTok{\# Get the class of \textquotesingle{}y\textquotesingle{}}
\FunctionTok{class}\NormalTok{(y)}
\end{Highlighting}
\end{Shaded}

\begin{verbatim}
[1] "numeric"
\end{verbatim}

\begin{Shaded}
\begin{Highlighting}[]
\CommentTok{\# Convert \textquotesingle{}y\textquotesingle{} from numeric to character type}
\NormalTok{y }\OtherTok{\textless{}{-}} \FunctionTok{as.character}\NormalTok{(y)}

\CommentTok{\# Get the class of \textquotesingle{}y\textquotesingle{} after conversion}
\FunctionTok{class}\NormalTok{(y)}
\end{Highlighting}
\end{Shaded}

\begin{verbatim}
[1] "character"
\end{verbatim}

\emph{Program code ends}

When we need to combine (concatenate) different strings to create longer ones, the
\texttt{paste()} function in R comes in handy. This has two
primary arguments: \texttt{sep} and \texttt{collapse}. The \texttt{sep}
argument specifies a separator to be placed between each pair of strings
being brought together. This can be an empty space, a dot, a dash,
or whichever symbol, word, or character you want to put between the
strings. On the other hand, when set to something other than its default
value of NULL, the collapse argument introduces an optional character
string to separate the results. For example, if we
bring together the string `My favorite book' with `has many chapters' using
\texttt{paste()}, we can create a seamless single string from these two.
In another scenario, we might use \texttt{paste()} to append the word
`Chapter' to numbers generated from 1 to 5, separating each resulting
string with a comma. Here is how it looks in R:

\emph{Program code starts}

\begin{Shaded}
\begin{Highlighting}[]
\CommentTok{\# Concatenating two strings}
\NormalTok{x }\OtherTok{\textless{}{-}}\FunctionTok{paste}\NormalTok{(}\StringTok{"My favorite book"}\NormalTok{, }\StringTok{"has many chapters"}\NormalTok{, }\AttributeTok{sep =} \StringTok{" "}\NormalTok{)}
\NormalTok{x}
\end{Highlighting}
\end{Shaded}

\begin{verbatim}
[1] "My favorite book has many chapters"
\end{verbatim}

\begin{Shaded}
\begin{Highlighting}[]
\CommentTok{\# Concatenating \textquotesingle{}Chapter\textquotesingle{} with numbers 1 to 5, separated by commas}
\NormalTok{x }\OtherTok{\textless{}{-}} \FunctionTok{paste}\NormalTok{(}\StringTok{"Chapter"}\NormalTok{, }\DecValTok{1}\SpecialCharTok{:}\DecValTok{5}\NormalTok{, }\AttributeTok{sep =} \StringTok{" "}\NormalTok{, }\AttributeTok{collapse =} \StringTok{", "}\NormalTok{)}
\NormalTok{x}
\end{Highlighting}
\end{Shaded}

\begin{verbatim}
[1] "Chapter 1, Chapter 2, Chapter 3, Chapter 4, Chapter 5"
\end{verbatim}

\emph{Program code ends}

You may have noticed that we used the same variable name, \texttt{x}, as
in the previous example. It is important to remember that reusing a
variable name in your code will overwrite its previous content with the
new data. Therefore, carefully tracking the names you assign to your
variables is good practice. Overwriting variables can lead to
confusion or errors, particularly in more complex scripts where the same
variable name might be used multiple times unintentionally. Thus, using
descriptive and unique names for your variables is advisable. This
approach helps to identify their purpose clearly and significantly
reduces the risk of accidental overwrites, thereby enhancing the
longevity and readability of your code.

\subsection{Working with strings}

In text mining, strings are our bread and butter, and to enhance our
understanding, in this section we will explore some simple examples of text
manipulation. This will involve using functions that
come with the base version of R, as well as additional functions that
can be installed to further enhance R's text-processing capabilities. We
will begin by importing the \textbf{Product Reviews} dataset, as
described in the Introduction, focusing on the `title' column.
To illustrate what the data looks like, we will display a small subset
(20 items). A larger dataset would be impractical to show in a printed
book. The data are selected randomly, but we ensure the reproducibility of
the results by using the \texttt{set.seed()} function. This approach
maintains a controlled randomness, ensuring consistent results across
different runs on different computers. Throughout this chapter, we will
demonstrate how the functions of the base version of R and the additional packages can
be used effectively in text mining to extract meaningful insights from our data.

\emph{Program code starts}

\begin{Shaded}
\begin{Highlighting}[]
\CommentTok{\# Set global option to not automatically convert strings to factors}
\FunctionTok{options}\NormalTok{(}\AttributeTok{stringsAsFactors =} \ConstantTok{FALSE}\NormalTok{)}

\CommentTok{\# Read the CSV file \textquotesingle{}Product Reviews.csv\textquotesingle{} into a dataframe \textquotesingle{}prodreviews\textquotesingle{}}
\NormalTok{prodreviews }\OtherTok{\textless{}{-}} \FunctionTok{read.csv}\NormalTok{(}\StringTok{"data/Product Reviews.csv"}\NormalTok{)}

\CommentTok{\# Convert the \textquotesingle{}Title\textquotesingle{} column of \textquotesingle{}prodreviews\textquotesingle{} to character type }
\CommentTok{\# and store in \textquotesingle{}prodtitle\textquotesingle{}}
\NormalTok{prodtitle }\OtherTok{\textless{}{-}} \FunctionTok{as.character}\NormalTok{(prodreviews}\SpecialCharTok{$}\NormalTok{Title)}

\CommentTok{\# Set a seed for random number generation to ensure reproducibility}
\FunctionTok{set.seed}\NormalTok{(}\DecValTok{123}\NormalTok{)}

\CommentTok{\# Sample 20 titles randomly from \textquotesingle{}prodtitle\textquotesingle{} and store them in \textquotesingle{}prodtitle20\textquotesingle{}}
\NormalTok{(prodtitle20 }\OtherTok{\textless{}{-}} \FunctionTok{sample}\NormalTok{(prodtitle, }\DecValTok{20}\NormalTok{))}
\end{Highlighting}
\end{Shaded}

\begin{verbatim}
 [1] "Best tech toy of 2010-2011"                              
 [2] "Worked for one day!"                                     
 [3] "Ignore the nonsense"                                     
 [4] "Nonna162"                                                
 [5] "Obviously only for the Tech Savvy"                       
 [6] "eh.."                                                    
 [7] "Excellent for what it is!"                               
 [8] "Does not add up at all"                                  
 [9] "Love the product!"                                       
[10] "Good purchase"                                           
[11] "Great Tablet"                                            
[12] "Not what I expected"                                     
[13] "Less than satisfaction"                                  
[14] "Boss tablet."                                            
[15] "Sweet Little Tablet"                                     
[16] "I have learned my lesson."                               
[17] "Lightweight 10-inch Android 2.3 tablet with so-so screen"
[18] "Kindle?"                                                 
[19] "The Kindle Fire is a great first effort"                 
[20] "Wifi version inferior to the 3g version"                 
\end{verbatim}

\emph{Program code ends}

As we can see, the system displays the titles in a reproducible
pseudo-random sequence, thanks to the use of \texttt{set.seed()}. It is
important to note that text in R is case-sensitive; a computer will consider 
the words `great' and `Great' to be two distinct entities. This
distinction is essential in text analysis and manipulation. To get a bit
of practice, let us use the \texttt{substr()} function to select a part
of these strings. For instance, our goal could be to extract the first
four letters of each title string. Technically, we achieve this by
specifying the start and end index positions in the \texttt{substr()}
function. In our case, we start from index position 1 (the first
letter of a word) and end at position 4.

\emph{Program code starts}

\begin{Shaded}
\begin{Highlighting}[]
\CommentTok{\# Extracting the first four letters of each title}
\FunctionTok{substr}\NormalTok{(prodtitle20, }\AttributeTok{start =} \DecValTok{1}\NormalTok{, }\AttributeTok{stop =} \DecValTok{4}\NormalTok{)}
\end{Highlighting}
\end{Shaded}

\begin{verbatim}
 [1] "Best" "Work" "Igno" "Nonn" "Obvi" "eh.." "Exce" "Does" "Love" "Good"
[11] "Grea" "Not " "Less" "Boss" "Swee" "I ha" "Ligh" "Kind" "The " "Wifi"
\end{verbatim}

\emph{Program code ends}

By abbreviating the titles in our dataset, we might encounter a
situation where some strings end up with identical sequences. These can be composed of a mix of lowercase or uppercase letters,
spaces, and punctuation marks. To mitigate the potential confusion
caused by such duplication, we can use the \texttt{abbreviate()}
function. This creates abbreviations from each string,
typically selecting the first letter of each word and ensuring a minimum
of four letters are used. If a string has fewer than four words,
\texttt{abbreviate()} automatically constructs abbreviations to match this
length.

Beyond reducing duplication, the \texttt{abbreviate()} function serves
several other purposes. It can save time, reduce computer memory usage,
and enhance data anonymity. It is also handy for creating concise labels
for graphs or managing large datasets where conserving space is critical.

\emph{Program code starts}

\begin{Shaded}
\begin{Highlighting}[]
\CommentTok{\# Apply the \textquotesingle{}abbreviate()\textquotesingle{} function to the \textquotesingle{}prodtitle20\textquotesingle{} vector}
\FunctionTok{abbreviate}\NormalTok{(prodtitle20)}
\end{Highlighting}
\end{Shaded}

\begin{verbatim}
                              Best tech toy of 2010-2011 
                                                 "Btto2" 
                                     Worked for one day! 
                                                  "Wfod" 
                                     Ignore the nonsense 
                                                  "Igtn" 
                                                Nonna162 
                                                  "N162" 
                       Obviously only for the Tech Savvy 
                                                "OoftTS" 
                                                    eh.. 
                                                  "eh.." 
                               Excellent for what it is! 
                                                 "Efwii" 
                                  Does not add up at all 
                                                "Dnauaa" 
                                       Love the product! 
                                                  "Ltp!" 
                                           Good purchase 
                                                  "Gdpr" 
                                            Great Tablet 
                                                  "GrtT" 
                                     Not what I expected 
                                                  "NwIe" 
                                  Less than satisfaction 
                                                  "Lsts" 
                                            Boss tablet. 
                                                  "Btb." 
                                     Sweet Little Tablet 
                                                  "SwLT" 
                               I have learned my lesson. 
                                                 "Ihlml" 
Lightweight 10-inch Android 2.3 tablet with so-so screen 
                                              "L1A2twss" 
                                                 Kindle? 
                                                  "Knd?" 
                 The Kindle Fire is a great first effort 
                                              "TKFiagfe" 
                 Wifi version inferior to the 3g version 
                                               "Wvitt3v" 
\end{verbatim}

\emph{Program code ends}

In this example, each title in the \texttt{prodtitle20} vector is
abbreviated, with each abbreviation having a minimum of four characters.

The length of strings in a dataset is a valuable statistical item,
particularly in text analysis. For instance, we might be curious about
the number of characters in each review title within our sample, or we
might want to identify the title with the highest character count. To
accomplish this, we can use the \texttt{nchar()} function to determine
the length of each string and the max function to find the maximum
length in the \texttt{prodtitle20} dataset. Applying these functions is
an important step for a marketing analyst examining text data: It helps them 
understand how extensively people are talking about a product in the
reviews. Longer titles might indicate more detailed discussions
or descriptions, which can be insightful for understanding customer
engagement and interest.

\emph{Program code starts}

\begin{Shaded}
\begin{Highlighting}[]
\CommentTok{\# Calculate the length of each title}
\FunctionTok{nchar}\NormalTok{(prodtitle20)}
\end{Highlighting}
\end{Shaded}

\begin{verbatim}
 [1] 26 19 19  8 33  4 25 22 17 13 12 19 22 12 19 25 56  7 39 39
\end{verbatim}

\begin{Shaded}
\begin{Highlighting}[]
\CommentTok{\# Find the maximum length among these titles}
\FunctionTok{max}\NormalTok{(}\FunctionTok{nchar}\NormalTok{(prodtitle20))}
\end{Highlighting}
\end{Shaded}

\begin{verbatim}
[1] 56
\end{verbatim}

\emph{Program code ends}

We obtained two results from our analysis: firstly, a vector of numbers
where each number represents the character count of each title in the
sample dataset; and, secondly, a single number indicating the maximum
length of the titles. The next step is determining which
title(s) correspond to this maximum length. For this, we use the
\texttt{which()} function in R, which identifies the position(s) of the
string(s) with the longest length. It is possible to find more than one
string sharing the same maximum number of characters, i.e., that of the
title of the longest string in the dataset. In our case, the longest
string has 56 characters.

\emph{Program code starts}

\begin{Shaded}
\begin{Highlighting}[]
\CommentTok{\# Extract the titles from \textquotesingle{}prodtitle20\textquotesingle{} that have the maximum length}
\NormalTok{prodtitle20[}\FunctionTok{which}\NormalTok{(}\FunctionTok{nchar}\NormalTok{(prodtitle20)}\SpecialCharTok{==}\FunctionTok{max}\NormalTok{(}\FunctionTok{nchar}\NormalTok{(prodtitle20)))]}
\end{Highlighting}
\end{Shaded}

\begin{verbatim}
[1] "Lightweight 10-inch Android 2.3 tablet with so-so screen"
\end{verbatim}

\emph{Program code ends}

The code can be translated from R into plain English as follows: \textbf{which} is
the string that has the \textbf{max}imum \textbf{n}umber of
\textbf{char}acters in the \textbf{prodtitle20} sample dataset?

We can also determine the total number of titles in the dataset and find out how
many are unique. To do these, we can utilize R's
\texttt{length()} and \texttt{unique()} functions. The \texttt{length()}
function will tell us the total number of titles in the dataset, while
\texttt{unique()} will help us identify the number of unique titles. We
apply these functions to the entire dataset using \texttt{prodtitle}.

\emph{Program code starts}

\begin{Shaded}
\begin{Highlighting}[]
\CommentTok{\# Count total titles in \textquotesingle{}prodtitle\textquotesingle{}}
\FunctionTok{length}\NormalTok{(prodtitle)}
\end{Highlighting}
\end{Shaded}

\begin{verbatim}
[1] 40741
\end{verbatim}

\begin{Shaded}
\begin{Highlighting}[]
\CommentTok{\# Count unique titles in \textquotesingle{}prodtitle\textquotesingle{}}
\FunctionTok{length}\NormalTok{(}\FunctionTok{unique}\NormalTok{(prodtitle))}
\end{Highlighting}
\end{Shaded}

\begin{verbatim}
[1] 31405
\end{verbatim}

\emph{Program code ends}

We can see a total of 40,741 titles in our dataset, out of which
31,405 are unique. This shows that the remaining 9,336 titles are
duplicates, i.e., they appear more than once.

Having mentioned the relevance of statistics in relation to this
dataset, we can gain a quick statistical overview using the
\texttt{summary()} function. This function is particularly useful for
analyzing the length of strings in the prodtitle dataset. By applying
\texttt{summary()} to the character lengths of the titles, we shift our
focus to a broader analysis of the entire dataset. This allows
us to understand not just the count of titles but also the distribution
of their lengths, providing insights into the dataset's characteristics.

\emph{Program code starts}

\begin{Shaded}
\begin{Highlighting}[]
\CommentTok{\# Get statistical summary of title lengths in \textquotesingle{}prodtitle\textquotesingle{}}
\FunctionTok{summary}\NormalTok{(}\FunctionTok{nchar}\NormalTok{(prodtitle))}
\end{Highlighting}
\end{Shaded}

\begin{verbatim}
   Min. 1st Qu.  Median    Mean 3rd Qu.    Max. 
   1.00   13.00   21.00   24.64   32.00  128.00 
\end{verbatim}

\emph{Program code ends}

We immediately notice that the lengths of strings in our dataset vary
significantly: the shortest title has only 1 character, while the
longest has 128 characters. The average (mean) number of
letters per title is 24.64, with a median of 21. This range in string
lengths could affect how we process and analyze our text data. For
instance, very short titles might need a closer look to ensure they are
meaningful, while the longer titles could be more descriptive or
complex, and rich in information.

But how are these string lengths distributed across our dataset?
And, how many strings are there of each length? To answer these
questions, let us plot a histogram, which is ideal for uncovering
patterns in datasets. Histograms are very helpful for showing the frequency of
groups, or `bins', of data. When creating our histogram, we will adjust
the default bin width by setting the \texttt{breaks} argument. This will
control the number of bins used to group our data. For our analysis, we
aim to distribute the data across 20 bins. However, it is important to
note that the number of breaks specified in R is treated as more of a
guideline than a strict parameter. R's internal algorithms may adjust
the number to produce a visually appealing histogram. While there are methods
of forcing R to produce an exact number of bins, as discussed in Ogle (2018), the
default behavior typically yields satisfactory results.

\emph{Program code starts}

\begin{Shaded}
\begin{Highlighting}[]
\CommentTok{\# Plot a histogram of string lengths in \textquotesingle{}prodtitle\textquotesingle{}}
\FunctionTok{hist}\NormalTok{(}\FunctionTok{nchar}\NormalTok{(prodtitle), }
     \AttributeTok{main =} \StringTok{\textquotesingle{}Histogram of Title Lengths\textquotesingle{}}\NormalTok{,  }\CommentTok{\# Set the main title of the histogram}
     \AttributeTok{breaks =} \DecValTok{20}\NormalTok{,                          }\CommentTok{\# Specify the number of bins as 20}
     \AttributeTok{xlab =} \StringTok{\textquotesingle{}Number of Characters\textquotesingle{}}\NormalTok{,        }\CommentTok{\# Label for the x{-}axis}
     \AttributeTok{ylab =} \StringTok{\textquotesingle{}Frequency of Titles\textquotesingle{}}\NormalTok{)         }\CommentTok{\# Label for the y{-}axis}
\end{Highlighting}
\end{Shaded}

\begin{figure}[H]

\centering{

\includegraphics{2.Chapter_files/figure-pdf/fig-title-length-plot-1.pdf}

}

\caption{\label{fig-title-length-plot}Length of the titles in the
Product Reviews dataset \newline{} Authors' own figure}

\end{figure}%

\emph{Program code ends}

We can conclude from the histogram (see
Figure~\ref{fig-title-length-plot}) that most of the review titles in
this dataset, some 8,237, have lengths of between 10 and 15
characters. Beyond this range, the frequency of string lengths gradually
decreases. This trend indicates that shorter titles are more common; the
longer the title, the less frequently it appears in the dataset.

\emph{Exercise starts}

\textbf{Exercise}: Try to find the longest title in the entire dataset
using R. Determine whether there is only one or
multiple titles of the maximum length. Hint: You can use the
\texttt{nchar()} function to calculate the length of each title and then
identify the title(s) with the maximum length. Once you find the longest
title(s), explore it/them further: What do they have in common, and what
might their length indicate about the content or nature of these titles?

\emph{Exercise ends}

\subsection{Regular expressions in R}\label{regular-expressions-in-r}

The previous sections served as a warm-up, using R-based procedures to
introduce you to the fascinating world of text analytics. Next, we will
explore the basics of regular expressions, a powerful tool for
searching, parsing, and manipulating text. All the examples in this
section are designed to work with the base R functions, meaning you do not
need to install any extra packages.

Regular expressions, often abbreviated as regex or regexp, are used not
only in search engines and by word processors, but are also integral to many
programming languages for detecting, locating, or replacing information.
This section will illustrate the basic concept and some elementary
functions of regular expressions in the R computing environment
(Sanchez 2013). For a more detailed exploration, see Friedl (2006)
or Forta (2018).

R supports two types of regular expression paradigms: POSIX and
Perl-like. The POSIX standard, an acronym for `Portable Operating System
Interface,' comprises a set of IEEE (Institute of Electrical and
Electronics Engineers) standards. These standards aim to provide a
consistent and portable interface for programming and user environments
across various operating systems, including but not limited to
Unix-based systems. The POSIX standard has two types of regular
expressions: basic regular expressions (BRE) and extended regular
expressions (ERE). R uses ERE by default in its regular
expression functions, such as \texttt{grep}, \texttt{grepl}, \texttt{sub}, and \texttt{gsub}.

Perl had a large community that has significantly contributed to the
development of a rich set of regex libraries and documentation, and it was renowned for its early and robust integration of regular expressions. This
makes Perl's regex particularly well-suited for complex pattern-matching
tasks. In R, Perl-compatible regular expressions can be used by
setting the perl = TRUE option, which allows for the use of advanced regex features
within R's regex functions.

For further reading and a deeper understanding of these topics, see the
the following: Lewine (1991); Cicekli et al. (2010); and Gagolewski
(2022).

\subsection{The foundations of regular
expressions}\label{foundations-of-regular-expressions}

A regular expression is a sequence of characters and metacharacters
which together form a search pattern. For instance, a simple pattern like
\textbf{{[}a-z{]}+} is designed to find any lowercase letter. This
pattern represents a concise string that tests for a specific pattern
capable of extracting or replacing substrings. Regular expressions can
search for a variety of elements, including letters (e.g., `a', `b',
`c'), numbers (e.g., `0', `1', `5'), or other characters (e.g., `=',
`-'). To the untrained eye, these patterns might look as random as the
keystrokes of a cat walking across a keyboard. If you have ever
witnessed such a feline typist, you will know exactly what we mean. Rest
assured, no cats were harmed in the writing of this book :)

To demystify these seemingly cryptic search patterns, consider an
example: that of finding all the email addresses in a corpus. An email address, such
as john.doe@someserver.com, can be interpreted as a user (john doe)
located at a domain (someserver.com). Our task is to create a regular
expression that captures all the email addresses in the text, regardless of
the domain, whether it is gmail, yahoo, outlook, or others. Constructing
a regular expression for this purpose might look challenging at first.
For instance, a basic pattern to match email addresses could be
something like:
\textbf{{[}a-zA-Z0-9.\_\%+-{]}+@{[}a-zA-Z0-9.-{]}+\textbackslash.{[}a-zA-Z{]}\{2,\}}.
This pattern is designed to match the general structure of email
addresses, but it can be pretty tricky to read for those new to regular
expressions. To aid in understanding and visualizing these patterns, we
can use tools provided by the \texttt{re2r} package (Wenfeng 2020). This offers a way to visually break down and interpret each part of a
regular expression. Figure~\ref{fig-regex-re2r} shows an example
of what this visualization looks like.

\emph{Program code starts}

\begin{Shaded}
\begin{Highlighting}[]
\CommentTok{\# Load re2r package}
\FunctionTok{library}\NormalTok{(re2r)}

\CommentTok{\# Visualize email regex pattern}
\FunctionTok{show\_regex}\NormalTok{(}\StringTok{"}\SpecialCharTok{\textbackslash{}\textbackslash{}}\StringTok{b[a{-}zA{-}Z0{-}9.\_\%+{-}]+@[a{-}zA{-}Z0{-}9.{-}]+}\SpecialCharTok{\textbackslash{}\textbackslash{}}\StringTok{.[a{-}zA{-}Z]\{2,4\}}\SpecialCharTok{\textbackslash{}\textbackslash{}}\StringTok{b"}\NormalTok{)}
\end{Highlighting}
\end{Shaded}

\begin{figure}[H]

\centering{

\includegraphics{2.Chapter_files/figure-pdf/fig-regex-re2r-1.pdf}

}

\caption{\label{fig-regex-re2r}Visualization of a Regular Expression for
an Email Address \newline Authors' own figure}

\end{figure}%

\emph{Program code ends}

Note that the \texttt{re2r} package has been removed from CRAN and may
contain bugs, such as displaying quantifiers like
\texttt{\{2,4\}} inaccurately. We have installed it from \texttt{GitHub} for
educational purposes. While it has potential as a tool, you should be
aware of its limitations and bugs. Another R regex enhancement tool is
the \texttt{regexplain} (Gadenbuie 2018) addin for RStudio, which is
installable from GitHub. It provides a user-friendly regex experience
for crafting and understanding regular expressions within RStudio. It
use both base R functions and the \texttt{stringr} package
(Wickham 2022b), on which there is more towards the end of this chapter.

We notice that the regular expression contains a sequence of characters
designed to capture one or more email addresses (Goyvaerts and Levithan
2012). Let us decode this in plain English: The expression starts and
ends with a word boundary \textbackslash b, encapsulating the email
address. The user name consists of a combination of characters from
a to z (or A to Z), numbers from 0 to 9, and special characters like a
dot . , an underscore \_, a percentage \%, a plus +, and a hyphen -, in
any order. This is followed by the @ symbol found in all email
addresses. After this, the server name is specified, which can include
letters (both upper and lower case), numbers, dots, and hyphens.
Finally, the domain's extension is indicated by a dot followed by 2 to 4
letters (e.g., .com, .org). It is important to note that while this
pattern captures many email formats, it may not encompass all,
especially more complex ones like second-level domains (e.g.,
john.doe@myserver.ac.at). Thus, while this is effective for most common email
addresses, the regex might not detect every possible valid email
format. Nevertheless, by adjusting and expanding the pattern, it is
possible to tailor the regex to match a wider variety of valid email
formats, accommodating different structures and domain specifications.

Let us now delve into the characters and metacharacters used in regular
expressions in a structured manner. Regular expressions are composed of
characters, literally symbols or letters, and metacharacters
which have special meanings in the context of regex. One of the key aspects
of regular expressions is the use of quantifiers, which allows us to
specify the number of repetitions of characters in a pattern.

\subsubsection{Quantifiers}\label{quantifiers}

Quantifiers, as detailed in Sanchez (2013), define how many times a
particular character or sequence of characters need to appear in a string
for a match to be found. The various quantifiers and their functions are
summarized in Table~\ref{tbl-quantifiers1}. These enable us
to construct flexible and powerful search patterns in regular
expressions.

\begin{longtable}[]{@{}ll@{}}
\caption{Quantifiers}\label{tbl-quantifiers1}\tabularnewline
\toprule\noalign{}
Quantifier & Description \\
\midrule\noalign{}
\endfirsthead
\toprule\noalign{}
Quantifier & Description \\
\midrule\noalign{}
\endhead
\bottomrule\noalign{}
\endlastfoot
\texttt{*} & Matches at least 0 times \\
\texttt{+} & Matches at least 1 time \\
\texttt{?} & Matches at most 1 time \\
\texttt{\{n\}} & Matches exactly n times \\
\texttt{\{n,\}} & Matches at least n times \\
\texttt{\{n,\ m\}} & Matches between n and m times \\
\end{longtable}

The quantifier in a regular expression is placed after the character or
group of characters whose occurrences (repetitions) it is meant to
refer to. For example, let us create a character vector in R containing
seven similar strings. These strings are designed to illustrate the
effect of applying various quantifiers and how they affect pattern
matching.

\emph{Program code starts}

\begin{Shaded}
\begin{Highlighting}[]
\CommentTok{\# Create a vector of sample strings with varying patterns}
\NormalTok{strings }\OtherTok{\textless{}{-}} \FunctionTok{c}\NormalTok{(}\StringTok{"x"}\NormalTok{, }\StringTok{"xy"}\NormalTok{, }\StringTok{"xz"}\NormalTok{, }\StringTok{"xzy"}\NormalTok{, }\StringTok{"xzzy"}\NormalTok{, }\StringTok{"xzzzy"}\NormalTok{, }\StringTok{"xzzzzy"}\NormalTok{, }\StringTok{"xzzzzyzy"}\NormalTok{)}
\end{Highlighting}
\end{Shaded}

\emph{Program code ends}

To identify patterns, such as finding all the sentences that contain a
specific sequence of characters, we can use the \texttt{grep()}
function. For instance, if we want to locate the sequence `xy' in our
previously created vector, we would use
\texttt{grep(\textquotesingle{}xy\textquotesingle{},\ strings)}. This
function returns the indices of the vector elements containing a
specific sequence, allowing us to identify and extract matching patterns
quickly.

\emph{Program code starts}

\begin{Shaded}
\begin{Highlighting}[]
\CommentTok{\# Identify indices of elements in vector \textquotesingle{}strings\textquotesingle{} that contain "xy"}
\FunctionTok{grep}\NormalTok{(}\StringTok{"xy"}\NormalTok{, strings)}
\end{Highlighting}
\end{Shaded}

\begin{verbatim}
[1] 2
\end{verbatim}

\emph{Program code ends}

The \texttt{grep()} function in R is used to identify patterns within a
vector of strings. By default, it returns the indices, i.e., the
positions in the vector, where the pattern you specify matches. If you
set the value parameter to \texttt{TRUE}, the function's behavior
changes: Instead of returning the indices, it will return the actual
strings that contain the pattern. This flexibility allows you to
either locate where the pattern occurs in the vector or extract
the matching elements.

\emph{Program code starts}

\begin{Shaded}
\begin{Highlighting}[]
\CommentTok{\# Return elements in \textquotesingle{}strings\textquotesingle{} that contain "xy"}
\FunctionTok{grep}\NormalTok{(}\StringTok{"xy"}\NormalTok{, strings, }\AttributeTok{value =} \ConstantTok{TRUE}\NormalTok{)}
\end{Highlighting}
\end{Shaded}

\begin{verbatim}
[1] "xy"
\end{verbatim}

\emph{Program code ends}

Let us analyze all the quantifying metacharacters and understand how
they perform in regular expressions.

The asterisk \texttt{*} matches no or more occurrences of the
preceding character. For instance, in a pattern, it can be used to match
any number of a specific character, such as z, that appears before
another character, like y. So, \texttt{z*y} in a regular expression
would match `y', `zy', `zzy', `zzzy', and so on, effectively capturing
sequences where `z' appears any number of times (including none) before a `y'.

\emph{Program code starts}

\begin{Shaded}
\begin{Highlighting}[]
\CommentTok{\# Return elements matching \textquotesingle{}xz*y\textquotesingle{} pattern in \textquotesingle{}strings\textquotesingle{}}
\FunctionTok{grep}\NormalTok{(}\StringTok{"xz*y"}\NormalTok{, strings, }\AttributeTok{value =} \ConstantTok{TRUE}\NormalTok{)}
\end{Highlighting}
\end{Shaded}

\begin{verbatim}
[1] "xy"       "xzy"      "xzzy"     "xzzzy"    "xzzzzy"   "xzzzzyzy"
\end{verbatim}

\emph{Program code ends}

The function successfully captured all the items in the vector where the
letter `y' is preceded by no or more occurrences of `z'.

Next, let us look at the plus character \texttt{+}. Unlike the
asterisk, the plus sign matches the preceding character one or more
times. It requires the preceding character to appear at least once,
making it useful for patterns where the character's presence is
mandatory but can repeat.

\emph{Program code starts}

\begin{Shaded}
\begin{Highlighting}[]
\CommentTok{\# Return elements matching \textquotesingle{}xz+y\textquotesingle{} pattern in \textquotesingle{}strings\textquotesingle{}}
\FunctionTok{grep}\NormalTok{(}\StringTok{"xz+y"}\NormalTok{, strings, }\AttributeTok{value =} \ConstantTok{TRUE}\NormalTok{)}
\end{Highlighting}
\end{Shaded}

\begin{verbatim}
[1] "xzy"      "xzzy"     "xzzzy"    "xzzzzy"   "xzzzzyzy"
\end{verbatim}

\emph{Program code ends}

The question-mark character \texttt{?} in regular expressions shows that
the preceding character appears at least once or not at all. This means it allows for the
presence of the character to be optional, indicating that the character
might appear once or never in the search pattern.

\emph{Program code starts}

\begin{Shaded}
\begin{Highlighting}[]
\CommentTok{\# Return elements matching \textquotesingle{}xz?y\textquotesingle{} pattern in \textquotesingle{}strings\textquotesingle{}}
\FunctionTok{grep}\NormalTok{(}\StringTok{"xz?y"}\NormalTok{, strings, }\AttributeTok{value =} \ConstantTok{TRUE}\NormalTok{)}
\end{Highlighting}
\end{Shaded}

\begin{verbatim}
[1] "xy"  "xzy"
\end{verbatim}

\emph{Program code ends}

The quantifier \texttt{\{n\}}, enclosed in curly brackets, is used in
regular expressions to match the precise number of occurrences – `n' – of the preceding
character. This allows for precise control over how many times a
character appears consecutively in the string being matched.

\emph{Program code starts}

\begin{Shaded}
\begin{Highlighting}[]
\CommentTok{\# Find \textquotesingle{}xz\{3\}y\textquotesingle{} (three \textquotesingle{}z\textquotesingle{}s before \textquotesingle{}y\textquotesingle{}) in \textquotesingle{}strings\textquotesingle{}}
\FunctionTok{grep}\NormalTok{(}\StringTok{"xz\{3\}y"}\NormalTok{, strings, }\AttributeTok{value =} \ConstantTok{TRUE}\NormalTok{)}
\end{Highlighting}
\end{Shaded}

\begin{verbatim}
[1] "xzzzy"
\end{verbatim}

\emph{Program code ends}

The quantifier \texttt{\{n,\}} in regular expressions, followed by a
comma, is used to match at least `n' occurrences of the preceding
character, setting `n' as the lower boundary. The absence of an upper boundary
in this quantifier means it can match `n' or more occurrences,
effectively extending to any (potentially infinite) number of
repetitions.

\emph{Program code starts}

\begin{Shaded}
\begin{Highlighting}[]
\CommentTok{\# Find \textquotesingle{}xz\{2,\}y\textquotesingle{} (at least two \textquotesingle{}z\textquotesingle{}s before \textquotesingle{}y\textquotesingle{}) in \textquotesingle{}strings\textquotesingle{}}
\FunctionTok{grep}\NormalTok{(}\StringTok{"xz\{2,\}y"}\NormalTok{, strings, }\AttributeTok{value =} \ConstantTok{TRUE}\NormalTok{)}
\end{Highlighting}
\end{Shaded}

\begin{verbatim}
[1] "xzzy"     "xzzzy"    "xzzzzy"   "xzzzzyzy"
\end{verbatim}

\emph{Program code ends}

In regular expressions, the quantifier \texttt{\{n,m\}} is used to match
a character that appears at least `n' times but not more than `m' times.
This allows us to specify a range for the number of occurrences of the
preceding character, with `n' as the lower limit and `m' as the upper
limit.

\emph{Program code starts}

\begin{Shaded}
\begin{Highlighting}[]
\CommentTok{\# Find \textquotesingle{}xz\{0,2\}y\textquotesingle{} (0 to 2 \textquotesingle{}z\textquotesingle{}s before \textquotesingle{}y\textquotesingle{}) in \textquotesingle{}strings\textquotesingle{}}
\FunctionTok{grep}\NormalTok{(}\StringTok{"xz\{0,2\}y"}\NormalTok{, strings, }\AttributeTok{value =} \ConstantTok{TRUE}\NormalTok{)}
\end{Highlighting}
\end{Shaded}

\begin{verbatim}
[1] "xy"   "xzy"  "xzzy"
\end{verbatim}

\emph{Program code ends}

In conclusion, this section has explored the versatile and powerful role
of quantifiers in regular expressions within R. We have seen how
quantifiers like \texttt{*}, \texttt{+}, \texttt{?}, \texttt{\{n\}},
\texttt{\{n,\}}, and \texttt{\{n,m\}} help us to maintain precise control over
pattern matching by specifying the number of occurrences of characters.

\emph{Exercise starts}

\textbf{Exercise}: Find all the string titles in the \texttt{prodtitle20}
variable which contain the pattern `ire'. Use the \texttt{grep()}
function in R for this purpose. Once you have identified these titles,
continue your exploration by applying the quantifiers we have discussed
above. Experiment with these quantifiers to discover a variety of
patterns within the \texttt{prodtitle20} variable.

\emph{Exercise ends}

\subsubsection{Position of a pattern inside a
string}\label{position-of-pattern-inside-a-string}

In the following section, we will focus on anchors and boundaries in
regular expressions, as detailed in sources like Bryan (2020). These
elements are crucial for specifying the position of a string within a
sequence of characters. Anchors denote a string's start and end, such as the circumflex (\^{}) for the former and the
dollar sign (\$) for the latter.
Boundaries, on the other hand, define specific positions within the
string. A summary of these anchors and boundaries, along with their
descriptions, is provided in Table~\ref{tbl-position1}.

\begin{longtable}[]{@{}ll@{}}
\caption{Position of a Pattern Inside a
String}\label{tbl-position1}\tabularnewline
\toprule\noalign{}
Anchors and Boundaries & Description \\
\midrule\noalign{}
\endfirsthead
\toprule\noalign{}
Anchors and Boundaries & Description \\
\midrule\noalign{}
\endhead
\bottomrule\noalign{}
\endlastfoot
\texttt{\^{}} & Matches the beginning of a string \\
\texttt{\$} & Matches the end of a string \\
\texttt{\textbackslash{}b} & Matches a word boundary \\
\texttt{\textbackslash{}B} & Matches a non-word boundary \\
\end{longtable}

We use an anchor in regular expressions to locate strings which start or
end with a specific word or character. Conversely, a boundary is used to
find a position between a word character (like letters and numbers) and
a non-word character (like spaces or punctuation). While the beginning
or end of a string could be considered a boundary, in regular
expressions, anchors and boundaries are distinct concepts with specific
uses. For example, if we want to capture all the sentences which start with
the character sequence `Ki', we would use the \texttt{\^{}} anchor
followed by `Ki'. The \texttt{\^{}} anchor explicitly matches a string's
start, so \texttt{\^{}Ki} would find any string that begins with `Ki'.

\emph{Program code starts}

\begin{Shaded}
\begin{Highlighting}[]
\CommentTok{\# Find elements in \textquotesingle{}prodtitle20\textquotesingle{} starting with "Ki"}
\FunctionTok{grep}\NormalTok{(}\StringTok{"\^{}Ki"}\NormalTok{, prodtitle20, }\AttributeTok{value =} \ConstantTok{TRUE}\NormalTok{)}
\end{Highlighting}
\end{Shaded}

\begin{verbatim}
[1] "Kindle?"
\end{verbatim}

\emph{Program code ends}

To capture phrases which end with a specific sequence of characters, such as
`et', we use the \texttt{\$} anchor. This metacharacter is placed after
the character sequence and ensures that the match occurs at the end of
the string. For example, using \texttt{et\$} would match any string that
concludes with `et'.

\emph{Program code starts}

\begin{Shaded}
\begin{Highlighting}[]
\CommentTok{\# Find elements in \textquotesingle{}prodtitle20\textquotesingle{} ending with "et"}
\FunctionTok{grep}\NormalTok{(}\StringTok{"et$"}\NormalTok{, prodtitle20, }\AttributeTok{value =} \ConstantTok{TRUE}\NormalTok{)}
\end{Highlighting}
\end{Shaded}

\begin{verbatim}
[1] "Great Tablet"        "Sweet Little Tablet"
\end{verbatim}

\emph{Program code ends}

A word boundary in regular expressions, denoted as
\texttt{\textbackslash{}\textbackslash{}b}, detects a position where a
word character (like a digit, a letter, or an underscore) is adjacent to
a non-word character (such as the beginning of a string, a tab, or a
space). In R, the backslash \texttt{\textbackslash{}} is a special
escape character. Therefore, when specifying a word boundary in a
regular expression within R, double backslashes must be used to escape
it properly, resulting in
\texttt{\textbackslash{}\textbackslash{}\textbackslash{}\textbackslash{}b}.
This ensures that the regular expression engine interprets it
as a word boundary.

\emph{Program code starts}

\begin{Shaded}
\begin{Highlighting}[]
\CommentTok{\# Find elements in \textquotesingle{}prodtitle20\textquotesingle{} with the word "Fire" as a whole word}
\FunctionTok{grep}\NormalTok{(}\StringTok{"}\SpecialCharTok{\textbackslash{}\textbackslash{}}\StringTok{bFire}\SpecialCharTok{\textbackslash{}\textbackslash{}}\StringTok{b"}\NormalTok{, prodtitle20, }\AttributeTok{value =} \ConstantTok{TRUE}\NormalTok{)}
\end{Highlighting}
\end{Shaded}

\begin{verbatim}
[1] "The Kindle Fire is a great first effort"
\end{verbatim}

\emph{Program code ends}

The sequence
\texttt{\textbackslash{}\textbackslash{}bFire\textbackslash{}\textbackslash{}b}
successfully matches the term `Fire' as a standalone word. However, if
`Fire' were in lowercase, this pattern would not match terms like
`firefly', `crossfire', or `unfired', where `fire' is a substring of a
larger word. The pattern will match sequences like `fire!me' because `!'
is not considered a word character. The word boundary
\texttt{\textbackslash{}\textbackslash{}b} can be used flexibly: Placing
\texttt{\textbackslash{}\textbackslash{}b} before a term (e.g.,
\texttt{\textbackslash{}\textbackslash{}bfire}) will match cases like
`firefly', and placing it after a term (e.g.,
\texttt{fire\textbackslash{}\textbackslash{}b}) will match cases like
`crossfire'. You can use the word boundary at either or both ends of the
search sequence.

\texttt{\textbackslash{}\textbackslash{}B} is the complement of
\texttt{\textbackslash{}\textbackslash{}b}. It matches positions where
the adjacent characters are either both word characters or both non-word
characters. For example, \texttt{\textbackslash{}\textbackslash{}B}
would identify a match within strings like `\$(=\%@' or in cases where word
characters surround a term without forming a standalone word. This makes
\texttt{\textbackslash{}\textbackslash{}B} useful for detecting patterns
which are not distinct words or embedded within other characters.

\emph{Program code starts}

\begin{Shaded}
\begin{Highlighting}[]
\CommentTok{\# Find elements in \textquotesingle{}prodtitle20\textquotesingle{} where "able" is part of another word}
\FunctionTok{grep}\NormalTok{(}\StringTok{"}\SpecialCharTok{\textbackslash{}\textbackslash{}}\StringTok{Bable}\SpecialCharTok{\textbackslash{}\textbackslash{}}\StringTok{B"}\NormalTok{, prodtitle20, }\AttributeTok{value =} \ConstantTok{TRUE}\NormalTok{)}
\end{Highlighting}
\end{Shaded}

\begin{verbatim}
[1] "Great Tablet"                                            
[2] "Boss tablet."                                            
[3] "Sweet Little Tablet"                                     
[4] "Lightweight 10-inch Android 2.3 tablet with so-so screen"
\end{verbatim}

\emph{Program code ends}

In the example above, the pattern
\texttt{\textbackslash{}\textbackslash{}Bable\textbackslash{}\textbackslash{}B}
is designed to detect the sequence `able' when it is part of a larger
word and surrounded by other word characters. It will not identify `able'
as a standalone word. However, it is important to note that this pattern will not match `able' when it is directly followed or preceded by
non-word characters, such as a full stop or a space. Essentially,
\texttt{\textbackslash{}\textbackslash{}B} requires both sides of
`able', or any other collection of word characters, to be word characters, showing where `able' is embedded
within a word but not at the beginning or end of a word.

\subsubsection{Common operators}\label{common-operators}

Operators in regular expressions, such as the dot (\texttt{.}) or
character lists denoted by three dots embraced by square brackets
(\texttt{{[}...{]})}, are essential when defining specific criteria for matches.
These operators are the building blocks of regular expressions, allowing
for the construction of complex search patterns.
Table~\ref{tbl-character1} presents some of the most common operators
used in regular expressions. For a comprehensive list and a detailed
explanation of these operators, refer to the works of Sanchez (2013) and
Forta (2018).

\begin{longtable}[]{@{}
  >{\raggedright\arraybackslash}p{(\columnwidth - 2\tabcolsep) * \real{0.1831}}
  >{\raggedright\arraybackslash}p{(\columnwidth - 2\tabcolsep) * \real{0.8169}}@{}}
\caption{Common Regular Expression Characters and Their
Functions}\label{tbl-character1}\tabularnewline
\toprule\noalign{}
\begin{minipage}[b]{\linewidth}\raggedright
Character
\end{minipage} & \begin{minipage}[b]{\linewidth}\raggedright
Description
\end{minipage} \\
\midrule\noalign{}
\endfirsthead
\toprule\noalign{}
\begin{minipage}[b]{\linewidth}\raggedright
Character
\end{minipage} & \begin{minipage}[b]{\linewidth}\raggedright
Description
\end{minipage} \\
\midrule\noalign{}
\endhead
\bottomrule\noalign{}
\endlastfoot
\texttt{.} & Matches any character except newline \\
\texttt{{[}...{]}} & Matches any one of the characters inside the
brackets. A dash \texttt{-} can specify a range, like
\texttt{{[}a-z{]}} \\
\texttt{{[}\^{}...{]}} & Matches any character not listed inside the
brackets \\
\texttt{\textbar{}} & The \texttt{OR} operator: matches patterns on
either side of \texttt{\textbar{}} \\
\texttt{(...)} & Groups characters. Used with quantifiers (\texttt{?},
\texttt{*}, \texttt{+}, \texttt{\{n\}}). Groups can be referenced with
\texttt{\textbackslash{}N}, where \texttt{N} is the group number \\
\texttt{\textbackslash{}} & Escapes a special character's meaning, like
\texttt{\textbackslash{}\$}, \texttt{\textbackslash{}*},
\texttt{\textbackslash{}+}, etc. In R, use double backslashes, e.g.,
\texttt{\textbackslash{}\textbackslash{}*} \\
\end{longtable}

Examine the following examples carefully to understand
how the different elements in regular expressions work. The single dot
(\texttt{.}) is a versatile operator in regular expressions. It matches
single characters (except for newline characters). For instance, in
the pattern \texttt{f.e}, the dot (\texttt{.}) will match any character
that appears between `f' and `e', such as `fle', `fee', or `fze'.

\emph{Program code starts}

\begin{Shaded}
\begin{Highlighting}[]
\CommentTok{\# Match \textquotesingle{}fe\textquotesingle{} followed by any character in \textquotesingle{}prodtitle20\textquotesingle{}}
\FunctionTok{grep}\NormalTok{(}\StringTok{"fe."}\NormalTok{, prodtitle20, }\AttributeTok{value =} \ConstantTok{TRUE}\NormalTok{)}
\end{Highlighting}
\end{Shaded}

\begin{verbatim}
[1] "Wifi version inferior to the 3g version"
\end{verbatim}

\emph{Program code ends}

The square brackets in regular expressions match any of the characters contained within them. When a hyphen is used between two letters it creates a
range, selecting all the characters that lie between them in (English) alphabetical
order. For example, \texttt{{[}a-i{]}} matches any letter from `a' to
`i': a, b, c, d, e, f, g, h, or i. In a pattern such as
\texttt{p{[}ru{]}}, `p' can be followed by either an `r' or a `u'. In a
pattern like \texttt{p{[}a-i{]}}, `p' can be followed by any letter from
`a' to `i'.

\emph{Program code starts}

\begin{Shaded}
\begin{Highlighting}[]
\CommentTok{\# Match \textquotesingle{}p\textquotesingle{} followed by \textquotesingle{}r\textquotesingle{} or \textquotesingle{}u\textquotesingle{} in \textquotesingle{}prodtitle20\textquotesingle{}}
\FunctionTok{grep}\NormalTok{(}\StringTok{"p[ru]"}\NormalTok{, prodtitle20, }\AttributeTok{value =} \ConstantTok{TRUE}\NormalTok{)}
\end{Highlighting}
\end{Shaded}

\begin{verbatim}
[1] "Love the product!" "Good purchase"    
\end{verbatim}

\begin{Shaded}
\begin{Highlighting}[]
\CommentTok{\# Match \textquotesingle{}p\textquotesingle{} followed by any letter from \textquotesingle{}a\textquotesingle{} to \textquotesingle{}i\textquotesingle{} in \textquotesingle{}prodtitle20\textquotesingle{}}
\FunctionTok{grep}\NormalTok{(}\StringTok{"p[a{-}i]"}\NormalTok{, prodtitle20, }\AttributeTok{value =} \ConstantTok{TRUE}\NormalTok{)}
\end{Highlighting}
\end{Shaded}

\begin{verbatim}
[1] "Not what I expected"
\end{verbatim}

\emph{Program code ends}

The caret symbol (\texttt{\^{}}) inside square brackets (not at the
beginning of the expression) inverts the character class, selecting
characters that are not listed inside the brackets. For example, in a
pattern like \texttt{f{[}\^{}i{]}}, the `f' must not be followed by an
`i': It matches any character following `f' except for `i'. In contrast,
a pattern like \texttt{fi} simply looks for the exact sequence `fi'.

\emph{Program code starts}

\begin{Shaded}
\begin{Highlighting}[]
\CommentTok{\# Match \textquotesingle{}f\textquotesingle{} not followed by \textquotesingle{}i\textquotesingle{} in \textquotesingle{}prodtitle20\textquotesingle{}}
\FunctionTok{grep}\NormalTok{(}\StringTok{"f[\^{}i]"}\NormalTok{, prodtitle20, }\AttributeTok{value =} \ConstantTok{TRUE}\NormalTok{)}
\end{Highlighting}
\end{Shaded}

\begin{verbatim}
[1] "Best tech toy of 2010-2011"             
[2] "Worked for one day!"                    
[3] "Obviously only for the Tech Savvy"      
[4] "Excellent for what it is!"              
[5] "Less than satisfaction"                 
[6] "The Kindle Fire is a great first effort"
[7] "Wifi version inferior to the 3g version"
\end{verbatim}

\begin{Shaded}
\begin{Highlighting}[]
\CommentTok{\# Match \textquotesingle{}fi\textquotesingle{} in \textquotesingle{}prodtitle20\textquotesingle{}}
\FunctionTok{grep}\NormalTok{(}\StringTok{"fi"}\NormalTok{, prodtitle20, }\AttributeTok{value =} \ConstantTok{TRUE}\NormalTok{)}
\end{Highlighting}
\end{Shaded}

\begin{verbatim}
[1] "The Kindle Fire is a great first effort"
[2] "Wifi version inferior to the 3g version"
\end{verbatim}

\emph{Program code ends}

In regular expressions, the \texttt{OR} operator, represented by the
vertical bar (\texttt{\textbar{}}), is used to match strings that
contain one of multiple specified patterns. For example, using
\texttt{er\textbar{}re} in a regular expression will match all strings
that contain either the sequence `er' or `re'.

\emph{Program code starts}

\begin{Shaded}
\begin{Highlighting}[]
\CommentTok{\# Match \textquotesingle{}er\textquotesingle{} or \textquotesingle{}re\textquotesingle{} in \textquotesingle{}prodtitle20\textquotesingle{}}
\FunctionTok{grep}\NormalTok{(}\StringTok{"er|re"}\NormalTok{, prodtitle20, }\AttributeTok{value =} \ConstantTok{TRUE}\NormalTok{)}
\end{Highlighting}
\end{Shaded}

\begin{verbatim}
[1] "Ignore the nonsense"                                     
[2] "Great Tablet"                                            
[3] "Lightweight 10-inch Android 2.3 tablet with so-so screen"
[4] "The Kindle Fire is a great first effort"                 
[5] "Wifi version inferior to the 3g version"                 
\end{verbatim}

\emph{Program code ends}

Capturing parentheses (), are often understood less well than other
operators. Quite simply, they allow us to treat multiple characters as a single unit. This can be particularly useful when applying quantifiers for extracting
specific parts of a string. In the pattern \texttt{(xy)+}, the parentheses group together `xy', so the \texttt{+} quantifier applies to
the sequence `xy', and not just 'x' or `y'. This pattern will also match strings
such as `xyxy, 'xyxyxy', and so on, where `xy' repeats one or more
times.

\emph{Program code starts}

\begin{Shaded}
\begin{Highlighting}[]
\CommentTok{\# Match one or more repetitions of \textquotesingle{}ire\textquotesingle{} in \textquotesingle{}prodtitle20\textquotesingle{}}
\FunctionTok{grep}\NormalTok{(}\StringTok{"(ire)+"}\NormalTok{, prodtitle20, }\AttributeTok{value =} \ConstantTok{TRUE}\NormalTok{)}
\end{Highlighting}
\end{Shaded}

\begin{verbatim}
[1] "The Kindle Fire is a great first effort"
\end{verbatim}

\emph{Program code ends}

Capturing parentheses can work wonders in combination with
\texttt{sub()} and \texttt{gsub()}, which stand for
\textbf{sub}stitution and \textbf{g}lobal \textbf{sub}stitution
respectively. The \texttt{sub()} function replaces the first occurrence
of a pattern, while \texttt{gsub()} replaces all occurrences. When using
capturing parentheses with these functions, the text within the
parentheses is the target for substitution. Capturing parentheses
are used to create captured groups. Capturing parentheses are
represented by the \texttt{(} and \texttt{)} symbols, and a
group is formed by the characters enclosed within these parentheses,
such as \texttt{(xyz)}. We reference captured groups in the replacement
string using double backslashes followed by a group number (e.g.,
\textbackslash\textbackslash1 for the first group,
\textbackslash\textbackslash2 for the second group, etc.), where the
number corresponds to the order of the group. For example, in the
pattern (xy) (yz), \textbackslash\textbackslash1 refers to the group,
(xy), and \textbackslash\textbackslash2, refers to the group (yz). This
allows us to specify which part of the pattern should be
modified or which should remain \textbf{unmodified}. By default, the
\texttt{gsub()} function processes each element of a character vector
and returns a vector of the same length, including modified and
unmodified elements. To focus on certain parts, we can first select
the elements which match our criteria using a function like
\texttt{grep()}. After isolating these elements, we then apply
\texttt{gsub()} to them. This allows us to view and work with
the data that matches our substitution pattern.

\emph{Program code starts}

\begin{Shaded}
\begin{Highlighting}[]
\CommentTok{\# Identify elements containing \textquotesingle{}Kindle Fire\textquotesingle{}}
\NormalTok{matches }\OtherTok{\textless{}{-}} \FunctionTok{grep}\NormalTok{(}\StringTok{"Kindle Fire"}\NormalTok{, prodtitle20, }\AttributeTok{value =} \ConstantTok{TRUE}\NormalTok{)}

\CommentTok{\# Apply gsub only to matched elements}
\NormalTok{modified }\OtherTok{\textless{}{-}} \FunctionTok{gsub}\NormalTok{(}\StringTok{"(Kindle) Fire"}\NormalTok{, }\StringTok{"}\SpecialCharTok{\textbackslash{}\textbackslash{}}\StringTok{1 Big Fire"}\NormalTok{, matches)}

\CommentTok{\# Display modified elements}
\NormalTok{modified}
\end{Highlighting}
\end{Shaded}

\begin{verbatim}
[1] "The Kindle Big Fire is a great first effort"
\end{verbatim}

\emph{Program code ends}

In our example, we searched for the sequence `Kindle Fire'. Using
\texttt{gsub()}, we kept the `Kindle' part unchanged while substituting
`Fire' for `Big Fire'. This was achieved by capturing groups and
back-referencing: the `Kindle' part was captured in a group and then
referenced in the replacement string to remain fixed. Another powerful
application of this function is to swap two terms in a string. By placing each term in a
separate group, we can use back-references in the replacement pattern to
reverse their positions.

\begin{Shaded}
\begin{Highlighting}[]
\CommentTok{\# Swap \textquotesingle{}Kindle\textquotesingle{} and \textquotesingle{}Fire\textquotesingle{}}
\FunctionTok{gsub}\NormalTok{(}\StringTok{"(Kindle) (Fire)"}\NormalTok{, }\StringTok{"}\SpecialCharTok{\textbackslash{}\textbackslash{}}\StringTok{2 }\SpecialCharTok{\textbackslash{}\textbackslash{}}\StringTok{1"}\NormalTok{, matches)  }
\end{Highlighting}
\end{Shaded}

\begin{verbatim}
[1] "The Fire Kindle is a great first effort"
\end{verbatim}

When we want to search for characters which are also metacharacters in
regular expressions, these characters must `escape' because they
have special meanings. For example, you might need to look for a dot (\texttt{.}) in a string. Remember that the dot
is also a metacharacter used to match any character, so you need to
'escape' it to transform it into a simple character. In most regular
expression engines, escaping is done with a single backslash. However,
in R, because the backslash itself is an escape character, you need to use two backslashes to represent a true
backslash. The first backslash escapes the second, making it a true
backslash, and the second backslash then escapes the dot
(\texttt{\textbackslash{}\textbackslash{}.}), treating it as a true
dot.

\emph{Program code starts}

\begin{Shaded}
\begin{Highlighting}[]
\CommentTok{\# Match literal dots in \textquotesingle{}prodtitle20\textquotesingle{}}
\FunctionTok{grep}\NormalTok{(}\StringTok{"}\SpecialCharTok{\textbackslash{}\textbackslash{}}\StringTok{."}\NormalTok{, prodtitle20, }\AttributeTok{value =} \ConstantTok{TRUE}\NormalTok{)}
\end{Highlighting}
\end{Shaded}

\begin{verbatim}
[1] "eh.."                                                    
[2] "Boss tablet."                                            
[3] "I have learned my lesson."                               
[4] "Lightweight 10-inch Android 2.3 tablet with so-so screen"
\end{verbatim}

\emph{Program code ends}

\emph{Exercise starts}

\textbf{Exercise 1}: Using the Product Reviews dataset, focus on the
Title column. Find all titles that contain the word `love' and end
with an exclamation mark `!'. Once you identified these, use regular expressions
to replace the exclamation mark at the end of these titles with a
question mark `?'. Employ back-referencing in your substitution.

\textbf{Exercise 2}: Use grouping parentheses to match
both ``can'' and ``can't''. Explore how the optional character operator
\texttt{?} can be combined with parentheses to handle variations in word
endings.

\emph{Exercise ends}

\subsubsection{Character classes}\label{character-classes}

Character classes in regular expressions are powerful tools used to
match specific sets of characters, such as letters, numbers, and more.
In the POSIX standard for regular expressions, these classes are placed
within square brackets and include additional colons and a key name,
following the structure \texttt{{[}:keyname:{]}}. For example,
\texttt{{[}:digit:{]}} represents all numerical digits, and
\texttt{{[}:alpha:{]}} matches all alphabetical characters.
Table~\ref{tbl-charclas1} summarizes some of the most widely used POSIX
character classes (Sanchez 2013; Forta 2018).

\begin{longtable}[]{@{}
  >{\raggedright\arraybackslash}p{(\columnwidth - 2\tabcolsep) * \real{0.1944}}
  >{\raggedright\arraybackslash}p{(\columnwidth - 2\tabcolsep) * \real{0.8056}}@{}}
\caption{POSIX Character Classes and Their
Functions}\label{tbl-charclas1}\tabularnewline
\toprule\noalign{}
\begin{minipage}[b]{\linewidth}\raggedright
Character
\end{minipage} & \begin{minipage}[b]{\linewidth}\raggedright
Description
\end{minipage} \\
\midrule\noalign{}
\endfirsthead
\toprule\noalign{}
\begin{minipage}[b]{\linewidth}\raggedright
Character
\end{minipage} & \begin{minipage}[b]{\linewidth}\raggedright
Description
\end{minipage} \\
\midrule\noalign{}
\endhead
\bottomrule\noalign{}
\endlastfoot
\texttt{{[}:lower:{]}} & Matches lower-case letters \\
\texttt{{[}:upper:{]}} & Matches upper-case letters \\
\texttt{{[}:digit:{]}} & Matches digits,
\texttt{0\ 1\ 2\ 3\ 4\ 5\ 6\ 7\ 8\ 9} \\
\texttt{{[}:alpha:{]}} & Matches alphabetic characters \\
\texttt{{[}:alnum:{]}} & Matches alphanumeric characters \\
\texttt{{[}:xdigit:{]}} & Matches hexadecimal digits,
\texttt{0-9\ A-F\ a-f} \\
\texttt{{[}:blank:{]}} & Matches blank characters, i.e., space and
tab \\
\texttt{{[}:space:{]}} & Matches space characters: tab, newline,
vertical tab, form feed, carriage return, space \\
\texttt{{[}:punct:{]}} & Matches punctuation characters \\
\texttt{{[}:graph:{]}} & Matches graphical (human-readable)
characters \\
\texttt{{[}:print:{]}} & Matches printable characters \\
\texttt{{[}:cntrl:{]}} & Matches control characters \\
\end{longtable}

Note that POSIX character classes, shown with colons and square
brackets \texttt{{[}:...:{]}}, must be used within another set of square
brackets. The correct pattern for matching a specific character
class, such as punctuation marks, would be \texttt{{[}{[}:punct:{]}{]}}.
This pattern will match any character that is classified as a
punctuation mark. Let us search for the titles that contain punctuation
marks in the \textbf{prodtitle20} dataset.

\emph{Program code starts}

\begin{Shaded}
\begin{Highlighting}[]
\CommentTok{\# Find titles with punctuation in \textquotesingle{}prodtitle20\textquotesingle{}}
\FunctionTok{grep}\NormalTok{(}\StringTok{"[[:punct:]]"}\NormalTok{, prodtitle20, }\AttributeTok{value =} \ConstantTok{TRUE}\NormalTok{)}
\end{Highlighting}
\end{Shaded}

\begin{verbatim}
[1] "Best tech toy of 2010-2011"                              
[2] "Worked for one day!"                                     
[3] "eh.."                                                    
[4] "Excellent for what it is!"                               
[5] "Love the product!"                                       
[6] "Boss tablet."                                            
[7] "I have learned my lesson."                               
[8] "Lightweight 10-inch Android 2.3 tablet with so-so screen"
[9] "Kindle?"                                                 
\end{verbatim}

\emph{Program code ends}

\emph{Exercise starts}

\textbf{Exercise}: Try writing regular expressions to identify titles
with specific character types, like numerical digits or uppercase letters.
Explore several patterns, such as those with punctuation, and apply
these to the entire dataset. If you are up for a challenge use the whole
\textbf{prodtitle} dataset.

\emph{Exercise ends}

\subsubsection{Sequences}\label{sequences}

Shorthand character classes in regular expressions offer a concise
notation system, similar to the more expansive character classes
discussed earlier (refer to Table~\ref{tbl-sequences1}). These shorthand
notations can often be used interchangeably with their longer character
class counterparts. The syntax involves a backslash \textbackslash{}
followed by a specific letter representing a particular set of
characters. Remember, in R, you need to escape the backslash with
another backslash. For a more comprehensive understanding of these
shorthand character classes and their usage, consult Sanchez (2013) and 
Forta (2018).

\begin{longtable}[]{@{}ll@{}}
\caption{Shorthand Character Classes in Regular
Expressions}\label{tbl-sequences1}\tabularnewline
\toprule\noalign{}
Regex Shorthand & Description \\
\midrule\noalign{}
\endfirsthead
\toprule\noalign{}
Regex Shorthand & Description \\
\midrule\noalign{}
\endhead
\bottomrule\noalign{}
\endlastfoot
\texttt{\textbackslash{}d} & Matches a digit character \\
\texttt{\textbackslash{}D} & Matches a non-digit character \\
\texttt{\textbackslash{}s} & Matches a space character \\
\texttt{\textbackslash{}S} & Matches a non-space character \\
\texttt{\textbackslash{}w} & Matches a word character
\texttt{{[}A-Za-z0-9{]}} \\
\texttt{\textbackslash{}W} & Matches a non-word character
\texttt{{[}\^{}A-Za-z0-9{]}} \\
\texttt{\textbackslash{}t} & Matches a tab character \\
\texttt{\textbackslash{}n} & Matches a new line character \\
\texttt{\textbackslash{}r} & Matches a carriage return \\
\end{longtable}

\emph{Exercise starts}

\textbf{Exercise}: Identify titles that contain specific patterns.
First, look for those which include exactly two consecutive digits. Then, search
for the ones which contain non-word characters (characters that are not
letters, digits, or underscores).

\emph{Exercise ends}

Important: You have seen three types of symbols in this chapter which
are borrowed from mathematics: parentheses (), brackets {[}{]}, and
braces \{\}, in other books referred to as round, square, and curly
brackets. The first ones are used for grouping, the second define
character classes, and the third are used to quantify specific limits.

\subsubsection{Regular expressions with the stringr
package}\label{regular-expressions-with-the-stringr-package}

The \texttt{stringr} package, part of Hadley Wickham's
\texttt{tidyverse} suite of tools, can be installed independently or
accessed as part of the entire \texttt{tidyverse} collection (refer to
the introduction for installation details; Wickham 2022b). All the main
functions in \texttt{stringr} start with the prefix \texttt{str\_} and
are designed for efficient string manipulation and the application of regular expressions. One key advantage of using \texttt{stringr} is its
standardization across the two regex standards implemented in R, which
means you do not need to specify \texttt{perl=TRUE/FALSE} for regex
operations (Bryan 2020). Examples of such functions are:
\texttt{str\_detect()} for finding patterns, \texttt{str\_subset()} for
extracting matching elements, and \texttt{str\_split\_fixed()} for
splitting strings based on patterns. It is worth mentioning that the
function \texttt{str\_view()} checks the representation of a string and
how a pattern matches.

\subsubsection{Regular expressions in search, and search and replace within
RStudio}\label{regular-expressions-in-search-search-and-replace-within-rstudio}

RStudio helps in using regular expression patterns to
search for or replace text in R, Rnw, Rmd, or other types of file. To
practice your regex skills in RStudio, you can create a regex sequence
and then enable the `Regex' and `Wrap' options. For instance, consider
the string ``Text Analytics for Digital Marketing''. Using the capturing
parentheses expression \texttt{(Text)(.*?)(i)} in the search box will
match the sequence ``Text Analyti''. This pattern uses three groups. The
first one captures ``Text'', the second captures any examples of no character or
least possible number, and the third captures ``i'' at its
first occurrence (the least possible). Similarly
\texttt{(ext)(.*?)(i)}, will match ``ext Analyti''. The addition of the
middle group (\texttt{.*?)} makes the this `lazy', meaning it
matches the shortest possible sequence available rather than the longest. The main
difference between greedy (standard in R) and lazy (enabled in RStudio)
matching, is that the former aims to match the longest possible sequence,
while the latter matches the shortest. The (\texttt{?}) is the only way to instruct the regex engine to stop at the shortest possible
match for a specified quantifier. For substitution, you can replace the
content of one of the groups. For example, to change ``Text Analytics
for Digital Marketing'' to ``My Text Analytics for Digital Marketing'':

\begin{itemize}
\tightlist
\item
  In the search box, enter \texttt{(Text)(.*)(Marketing)}
\item
  In the replace box, enter
  \texttt{My\ Text\textbackslash{}2\textbackslash{}3}
\end{itemize}

This means that the second and third groups remain unchanged, and
only ``Text'' is replaced with ``My text''. While R's regex engines
(POSIX and Perl-like) do not inherently support lazy matching, RStudio,
as an integrated development environment, provides this functionality,
demonstrating the flexibility different tools offer in handling regular
expressions. As a short exercise using the lazy version, you can try to
substitute a string inside a sentence.

\emph{Exercise starts}

\textbf{Exercises}

Copy and paste the following code into RStudio to begin working with a
set of strings designed for practicing regular expressions:

\emph{Program code starts}

\begin{Shaded}
\begin{Highlighting}[]
\NormalTok{text\_strings }\OtherTok{\textless{}{-}} \FunctionTok{c}\NormalTok{(}\StringTok{"We have to extract these numbers 12, 47, 48"}\NormalTok{,}
\StringTok{"The integers numbers are also interestings: 189 2036 314"}\NormalTok{,}
\StringTok{"\textquotesingle{},\textquotesingle{}is a separator, so please extract these numbers 125,789,1450}
\StringTok{and also these 564,90456"}\NormalTok{, }\StringTok{"We like to to offer you 7890$ per}
\StringTok{month in order to complete this task... we are joking"}\NormalTok{, }\StringTok{"You are}
\StringTok{going to learn 3 things, the first one is not to extract, and}
\StringTok{2 and 3 are simply digits."}\NormalTok{, }\StringTok{"Have fun with our mighty test, you}
\StringTok{are going to support science, progress, mankind wellness and you}
\StringTok{are going to waste 30 or 60 minutes of your life."}\NormalTok{, }\StringTok{"you can also}
\StringTok{extract exotic stuff like a456 gb67 and 45678911ghth"}\NormalTok{, }\StringTok{"Writing 1}
\StringTok{example is not funny, please consider that 66\% is}
\StringTok{validation+testing"}\NormalTok{, }\StringTok{"You you are a genius, I think that you like}
\StringTok{arrays A LOT, [3,45,67,900,1974]"}\NormalTok{, }\StringTok{"Who loves arrays more than}
\StringTok{me?"}\NormalTok{, }\StringTok{"\{366,78,90,5\}Yes, there are only 4 numbers inside"}\NormalTok{,}
\StringTok{"Integers are fine but sometimes you like 99 cents after the 99}
\StringTok{dollars"}\NormalTok{, }\StringTok{"100 euro are better than 99 euro"}\NormalTok{, }\StringTok{"I like to give you}
\StringTok{1000 numbers now: 12 3 56 21 67, and more, [45,67,7]"}\NormalTok{, }\StringTok{"Ok ok 1 2}
\StringTok{3 4 5 and the last one is 6"}\NormalTok{, }\StringTok{"33 trentini entrarono a Trento,}
\StringTok{tutti e 33 di tratto in tratto trotterellando"}\NormalTok{)}
\end{Highlighting}
\end{Shaded}

\emph{Program code ends}

\textbf{Selection}

\begin{enumerate}
\def\labelenumi{\arabic{enumi}.}
\tightlist
\item
  Select all the strings with words containing exactly two consecutive
  't's.
\item
  Select all the strings with words containing exactly two 't's (not
  necessarily consecutive) or exactly two 'e's (not necessarily
  consecutive).
\item
  Select all the strings that end in exactly one digit.
\item
  Select all the strings that end in two digits.
\item
  Select all the strings that have more than one capital letter.
\item
  Select all the strings that do not have the substring `nu'.
\item
  Select all the strings that contain special characters such as {[}, €,
  and so on.
\item
  Select all the strings that do not contain special characters.
\item
  Select all the strings that end with a punctuation mark.
\item
  Select all the strings that start with a lowercase letter.
\end{enumerate}

\textbf{Substitution or addition}

\begin{enumerate}
\def\labelenumi{\arabic{enumi}.}
\tightlist
\item
  Add a punctuation mark to strings that end without one.
\item
  Replace ``dollars'' with the \$ symbol.
\item
  Replace ``euros'' with the € symbol.
\item
  Replace consecutively repeating words with a single occurrence (e.g.,
  ``Bye Bye'' becomes ``Bye'').
\item
  Replace digits separated by a space with digits separated by a comma.
\item
  Replace digits separated by a comma with digits separated by a comma
  and a space.
\item
  Capitalize words that start with a lowercase letter.
\item
  Add a dollar sign at the end of strings that finish with a dot.
\item
  Insert a caret sign before the dot at the end of strings.
\item
  Place a caret sign at the beginning of strings that start with a
  lowercase letter.
\end{enumerate}

Bonus Exercise: Identify and correct a phrase in the dataset that
contains clear grammatical errors using regular expressions.

\emph{Exercise ends}

\textbf{Conclusion}

\emph{Conclusion starts}

This chapter was not intended to be an exhaustive exploration of regular
expressions. Instead, it was designed to guide novice and intermediate
users through the process of importing and manipulating text data in R,
equipping you with fundamental and robust tools for finding, adding,
removing, or substituting characters in your texts. We have presented regular using examples and exercises aimed at
simplifying complex concepts to help you understand them easily. Additionally, we have
introduced some basic statistics related to text strings, illustrated
with the help of a dataset. With these skills and knowledge, you are now
well-prepared to move on to the next chapters. We will continue the
marketing analytics journey through this book by exploring essential
areas such as obtaining text and text preprocessing. These chapters will
build upon what you have learned so far.

\emph{Conclusion ends}